\section{Pinned Adaptive Aggregation}
\label{sec:optim-adaptive-tiled}

The second \arcto{} optimization targets the host-device transfer
path. The pageable baseline is simple but slow: on a 1\,GiB TTI
input on the RX 7900 XT (gfx1100), pageable Cascaded spends
154.98\,ms in H2D and 151.66\,ms in D2H, with a total of 324.15\,ms;
the kernel itself contributes only a few milliseconds. Snappy and LZ4
on the same input show the same pattern, with pageable totals of
346.28\,ms and 447.34\,ms respectively. Pinned host staging is the
standard remedy because it removes the page-pinning step from the
critical path and unlocks bulk PCIe bandwidth, but the question of
\emph{how much} pinned memory to allocate, and \emph{when}, controls
whether the gain materializes.

\subsection{Pinned full versus pinned bounded}
\label{sec:pinned-tradeoff}

The most direct way to use pinned memory is to allocate a host buffer
the same size as the input, copy the pageable source into it, and
issue a single bulk H2D. For the 1\,GiB TTI workload above, this
strategy brings H2D and D2H down to roughly 39\,ms each across the
three codecs and reduces the measured pipeline time from 324\,ms to
96\,ms for Cascaded, from 346\,ms to 118\,ms for Snappy, and from
447\,ms to 220\,ms for LZ4. The kernel-visible behavior is unchanged;
all the difference comes from the transfer path.

What the kernel-visible time hides is the cost of the pinned buffer
itself. The same Cascaded run pays 227.49\,ms of pinned allocation
and 96.55\,ms of host-to-host staging to copy the 1\,GiB input into
the pinned region. Snappy and LZ4 show the same overhead structure
(229--221\,ms of allocation and 96\,ms of staging). Whether these
costs matter depends on usage: if the pinned buffer is reused across
many compression calls, the allocation amortizes; if it is created
per call, the pinned advantage is more apparent than real. Either
way, the pinned-memory \emph{footprint} scales with the input. A
4\,GiB input requires a 4\,GiB pinned allocation, which is a scarce
host resource that may conflict with other components of the
application.

\arcto{} resolves this with an adaptive pinned aggregation path that
allocates a fixed-size pinned window once and reuses it across
successive portions of the input. The kernel-visible chunks
$C_{\text{ker}}$ remain small, as established in
Section~\ref{sec:optim-chunk-size}; many such chunks are grouped into
a single pinned window of size $W_{\text{agg}}$. For an input of
length $L$, the runtime processes
$N_{\text{windows}} = \lceil L / W_{\text{agg}} \rceil$ windows in
sequence, reusing the same pinned buffer. The peak pinned-memory
footprint is therefore $W_{\text{agg}}$ rather than $L$, regardless
of how large the input is.

\subsection{Three-floor rule for the aggregation window}
\label{sec:three-floor-rule}

The aggregation window cannot be chosen by a single criterion because
three distinct overheads must be amortized: kernel underutilization,
PCIe setup, and per-window launch bookkeeping. \arcto{} therefore
selects $W_{\text{agg}}$ as the maximum of three architecture-aware
floors, so that whichever overhead is binding is the one that sets
the window size:
\[
W_{\text{agg}}^{*} =
\max\left(W_{\text{kernel\_sat}},\ W_{\text{pcie\_amort}},\
W_{\text{launch\_amort}}\right).
\]
The kernel-saturation floor combines the kernel-visible chunk size
with the resident wave-slot count of the GPU,
$W_{\text{kernel\_sat}} = S_{\text{wave}} \cdot C_{\text{ker}}^{*}$,
so that each window contains at least one chunk per resident wave
slot. The PCIe-amortization floor, set to 64\,MiB, ensures that the
bulk H2D issued per window is large enough to reach the asymptotic
PCIe Gen4 bandwidth and amortize the per-transfer setup. The
launch-amortization floor, set to 16\,MiB, ensures that the
per-window kernel launch and bookkeeping overhead is small relative
to the per-window kernel time. The values used by the four AMD parts
in the evaluation are summarized in Table~\ref{tab:three-floors}.

\begin{table}[t]
\centering
\caption{Three-floor rule applied to the four AMD GPUs evaluated.
$S_{\text{wave}}$ is the resident wave-slot count;
$C_{\text{ker}}^{*}$ is the chunk size selected by
Section~\ref{sec:optim-chunk-size}; the binding floor is shown in
bold.}
\label{tab:three-floors}
\small
\begin{tabular}{lrrrrrr}
\toprule
GPU (arch) & $S_{\text{wave}}$ & $C_{\text{ker}}^{*}$
  & $W_{\text{kernel\_sat}}$ & $W_{\text{pcie\_amort}}$
  & $W_{\text{launch\_amort}}$ & $W_{\text{agg}}^{*}$ \\
\midrule
MI50 (gfx906)     &  960 & 16\,KiB & 15\,MiB &
                    \textbf{64\,MiB} & 16\,MiB & 64\,MiB \\
MI210 (gfx90a)    & 3328 & 16\,KiB & 52\,MiB &
                    \textbf{64\,MiB} & 16\,MiB & 64\,MiB \\
MI300X (gfx942)   & 9728 &  8\,KiB & \textbf{76\,MiB} &
                    64\,MiB & 16\,MiB & 76\,MiB \\
RX 7900 XT (gfx1100) & 2688 & 16\,KiB & 42\,MiB &
                    \textbf{64\,MiB} & 16\,MiB & 64\,MiB \\
\bottomrule
\end{tabular}
\end{table}

Two observations follow from the table. First, MI300X is the only
case in which $W_{\text{kernel\_sat}}$ is the binding constraint: its
large $S_{\text{wave}}$ pushes the kernel-saturation floor above the
64\,MiB PCIe floor, so the adaptive runtime grows the window to
76\,MiB to provide one 8\,KiB chunk per resident wave slot. Second,
on the three smaller AMD parts the PCIe-amortization floor wins, and
all three converge on a 64\,MiB window. This is a useful property of
the rule: the same runtime policy adapts to a much smaller wave
budget without producing a window so small that PCIe setup
dominates.

\subsection{Adaptive execution and measured impact}
\label{sec:adaptive-execution}

For each window, \arcto{} copies the next portion of the pageable
input into the pinned aggregation region, transfers the populated
window to device memory, launches the batched compression kernel over
the chunks contained in that window, and copies the compressed
payload and metadata back to the host. The same pinned buffer is
reused for the next window, so the allocation cost is paid once for
the lifetime of the adaptive context. This matters because pinned
allocation is the dominant fixed cost on AMD platforms: a 1\,GiB
pinned region takes more than 200\,ms to allocate, while a 64\,MiB
region takes around 15\,ms.

The behavior on the 1\,GiB TTI workload, RX 7900 XT (gfx1100), is
summarized in Table~\ref{tab:adaptive-impact}. Adaptive execution
recovers the H2D/D2H gain of the pinned-full path, matches its
measured pipeline time, and reduces the pinned-memory footprint from
the full input size to a single 64\,MiB window. The 1\,GiB input is
processed in 16 windows. Compared to pageable execution, the
end-to-end speedup is $2.03\times$ for LZ4, $2.93\times$ for Snappy,
and $3.38\times$ for Cascaded; compared to the pinned-full path,
adaptive uses $16\times$ less peak pinned memory at the same total
time. The pattern extends to larger inputs without changing the
runtime policy: a 4\,GiB input is processed in 64 windows on the
same GPU, and the peak pinned footprint stays at 64\,MiB.

\begin{table}[t]
\centering
\caption{Pageable, pinned-full, and adaptive execution on a 1\,GiB
TTI input on the RX 7900 XT (gfx1100). Total time is the measured
end-to-end pipeline (compression). Peak pinned is the host pinned
allocation maintained for the duration of the call.}
\label{tab:adaptive-impact}
\small
\begin{tabular}{llrrrr}
\toprule
Algorithm & Mode & H2D & D2H & Total & Peak pinned \\
          &      & (ms) & (ms) & (ms) & (MiB) \\
\midrule
LZ4      & pageable     & 155.0 & 149.8 & 447.3 & 0 \\
         & pinned full  &  39.9 &  38.6 & 220.3 & 1024 \\
         & adaptive     &  39.8 &  38.5 & 220.2 & 64 \\
\midrule
Snappy   & pageable     & 155.2 & 150.5 & 346.3 & 0 \\
         & pinned full  &  39.6 &  38.4 & 118.4 & 1024 \\
         & adaptive     &  39.3 &  38.1 & 118.1 & 64 \\
\midrule
Cascaded & pageable     & 155.0 & 151.7 & 324.2 & 0 \\
         & pinned full  &  39.7 &  38.9 &  96.1 & 1024 \\
         & adaptive     &  39.4 &  38.6 &  95.4 & 64 \\
\bottomrule
\end{tabular}
\end{table}

The adaptive path is exposed to applications through a small context
API. The context object owns the pinned window, the runtime
counters, and the iteration state, and is reused across many
compression calls; this is the common case in scientific workflows
that compress successive checkpoints, timesteps, or partitions of a
larger dataset. The same context emits the per-call statistics used
in the evaluation: kernel chunk size, aggregation window size,
number of windows, peak pinned bytes, allocation time,
host-to-host copy time, H2D and D2H times, kernel time, and total
time. Reporting this full breakdown is what makes the comparison in
Table~\ref{tab:adaptive-impact} honest: omitting allocation and
host-to-host staging, as a kernel-only or transfer-only view would,
hides exactly the cost that the adaptive path eliminates.
